\documentclass[10pt,a4paper]{amsart}
\usepackage[margin=19mm]{geometry}
\usepackage{amsmath,amssymb,mathtools}
\usepackage[scr=rsfs,cal=euler]{mathalfa}
\usepackage{xfrac}
\usepackage[unicode,hidelinks,bookmarks=false]{hyperref}
\newcommand{\ie}{i.e.\ }
\newcommand{\cf}{cf.\ }
\newcommand{\resp}{respectively\ }
\newcommand{\Subsection}{\subsection}
\providecommand{\nobreakdash}{\nobreak\hbox{-}}
\setlength{\emergencystretch}{2em}
\multlinegap0pt
\usepackage{amssymb}
\usepackage{xcolor}
\usepackage{tcolorbox}
\usepackage{mathtools}
\usepackage{tikz}
\newtheorem{proposition}{Proposition}
\def\A{\mathcal{A}}
\renewcommand{\L}[0]{\mathcal{L}}
\def\M{\mathcal{M}}
\def\NN{\mathbb{N}}
\def\P{\mathcal{P}}
\def\U{\mathcal{U}}
\title{The reverse mathematics of the pigeonhole hierarchy}
\author{Quentin Le Hou\'erou, Ludovic Levy Patey, Ahmed Mimouni}
\date{Source: arXiv:2407.01236v2 (CC BY 4.0)}
\begin{document}
\maketitle

\noindent\emph{Source and licence.} The reverse mathematics of the pigeonhole hierarchy. Version arXiv:2407.01236v2 (CC BY 4.0), distributed by arXiv under the Creative Commons Attribution 4.0 International licence, http://creativecommons.org/licenses/by/4.0/, as recorded in the arXiv licence metadata for this exact version and shown on the abstract page https://arxiv.org/abs/2407.01236v2. Full licence notice: \texttt{licenses/arxiv-2407.01236-CC-BY-4.0.txt}.\par\smallskip

\noindent\textit{Editorial scope.} Counted unit: the article's proposition proposition (source0.tex lines 863--866). The article's complete original proof of it is delivered with it (source0.tex lines 868--908, one \texttt{proof} environment of 41 lines). No mathematics is written, repaired or re-derived: the statement and the proof are the article's own lines, in the article's own notation, and no retained line is substituted or re-flowed.  No further source-local context is needed: every cross-reference of every retained line resolves inside the two delivered spans.  Nothing was omitted from any retained span, so the package carries no omission record.  No retained line contains a citation, so the wrapper carries no bibliography and no bibliography entry.  No retained line contains a figure environment, an image-inclusion command or drawing code, so no image is delivered and the package declares no asset dependency.  Editorial formatting only, in addition: an AMS \texttt{amsart} preamble that restates the article's own declarations and macros used by the retained lines, namely the environments \texttt{proposition} on separate counters, together with the standard packages those lines need.  The retained environments print in this document as Proposition 1 in order of appearance, whereas the article prints the same items with the numbering of its own full document.  The complete native source of the article as distributed in the arXiv e-print for this exact version is bundled unchanged as \texttt{sources/161628/source0.tex}.
\begin{proposition}
For every countable Turing ideal~$\M$, there exists a set $C$ such that $\U_C^\M$ is $\M$-cohesive but not $\M$-minimal.
%    There exists an $\M$-cohesive class (of the form $\U_C^{\M}$) that is not $\M$-minimal.
\end{proposition}

\begin{proof}
Consider the following $\Pi_2^0$ class $\P = \{X : (\forall n)(\exists x) |[x,2^x) \cap X| \geq n\}$.
\smallskip

\emph{Claim 1. The class $\P$ is partition regular.} First, $\NN \in \P$. Let~$X_0 \cup X_1 \in \P$ for some sets $X_0,X_1$, if $X_0,X_1 \notin \P$, then for all $i < 2$, there exists some $n_i \in \NN$ such that for all $x$, and  $|[x,2^x) \cap X_i| < n_i$. Thus, for all $x \in \NN$, $|[x,2^x) \cap (X_0 \cup X_1)| < n_0 + n_1$, contradicting our assumption that $X_0 \cup X_1 \in \P$. This proves our claim.
\smallskip

Let $\U_D^\M$ be an $\M$-cohesive partition regular subclass of~$\P$, which exists as $\P$ is $\Pi^0_2$.
Let $Z_0, Z_1, \dots$ be the list of all sets~$X$ in $\M$ such that $\U_D^\M \subseteq \L_X$ and
let $\U_C^\M = \bigcap_n \L_{Z_n}$.
By $\M$-cohesiveness of $\U_D^\M$, for every $X \in \M$, either $\U_D^\M \subseteq \L_X$ or $\U_D^\M \subseteq \L_{\overline{X}}$, so $\U_C^\M$ is $\M$-cohesive. Moreover, $\U_D^\M \subseteq \U_C^\M$.
\smallskip

\emph{Claim 2. $\U_C^\M \not \subseteq \U_D^\M$}
Let $X = \{x_0, x_1, \dots \}$ defined as follows: let $x_n$ be the smallest element of $\bigcap_{k \leq n} Z_k$ bigger than $2^{x_{n-1}}$ (or bigger than $0$ if $n = 0$). Note that $\bigcap_{k \leq n} Z_k$ is infinite, since $\{ \overline{Z}_k : k \leq n \} \cup \{\bigcap_{k \leq n} Z_k\}$ is a cover of~$\NN$ and none of $\overline{Z}_k$ belongs to $\U_C^\M$, so $\bigcap_{k \leq n} Z_k \in \U_C^\M \subseteq \L_\NN$. By construction $X$ cannot be an element of $\P$ as $|X \cap [x,2^x)| \leq 1$ for every $x$, and $X \in \U_C^\M$ as $X \cap Z_k$ is infinite for every $k \in \NN$.
It follows that $\U_C^\M$ is $\M$-cohesive, but not $\M$-minimal.


% Say $\M = \{Z_0, Z_1, \dots\}$. For every $n \in \NN$ and $i < 2$, let $Z_n^0 = Z_n$ and $Z_n^1 = \overline{Z}_n$.
% Let $\sigma_0 \prec \sigma_1 \prec \dots$ be an infinite sequence of strings with $|\sigma_k| = k$ defined inductively as follows. First, $\sigma_0 = \epsilon$. Assume $\sigma_k$ has been defined so that $\bigcap_{s < k} Z_s^{\sigma_k(s)} \in \P$. Then by partition regularity of~$\P$, there is some~$i < 2$ such that $Z_k^i \cap \bigcap_{s < k} Z_s^{\sigma_k(s)} \in \P$. Let $\sigma_{k+1} = \sigma_k \cdot i$. Let $H$ be the limit of $\sigma_0, \sigma_1, \dots$ and let $\A = \bigcap_{i \in \NN} \L_{Z_i^{H(i)}}$.
% \smallskip

% \emph{Claim 2. $\A \cap \M = \{ Z_k^{H(k)} : k \in \NN \}$.} Fix~$k \in \NN$. We first prove that $Z_k^{H(k)} \in \A$. Assume otherwise, then there exists some $\ell$ such that $Z_k^{H(k)} \cap Z_\ell^{H(\ell)}$ is finite, contradicting the fact that $\bigcap_{s \leq \max\{k,\ell\}} Z_s^{H(s)} \in \P$. It is also clear that $Z_k^{1 - H(k)} \notin \A$, since $\A \subseteq \L_{Z_k^{H(k)}}$. This proves our claim. It follows that $\A \cap \M \subseteq \P$.
% \smallskip

% \emph{Claim 3. $\A \cap \M$ is closed under intersection.} By Claim 2, let $Y_0 = Z_{i_0}^{H(i_0)}$ and $Y_1 = Z_{i_1}^{H(i_1)}$ be two elements of $\A \cap \M$. As $\M$ is a Turing ideal, there exists some index $i_2$ such that $Z_{i_2} = Y_0 \cap Y_1$. If $H(i_2) = 1$, then, by construction of $\A$, $Y_0 \cap Y_1 \cap \overline{Y_0 \cap Y_1} = \emptyset$ would be an element of $\P$, which is not the case. Hence, $H(i_2) = 0$ and $Y_0 \cap Y_1 \in \A \cap \M$. 
% \smallskip


% \emph{Claim 4. $\A$ is partition regular.} First, $\NN \in \A$. Assume by contradiction that there exist some sets $X_0,X_1$ such that $X_0 \cup X_1 \in \A$ and $X_0, X_1 \notin \A$. By definition of $\A$, there exist some indices $i_0,i_1$ such that $X_0 \cap Z_{i_0}^{H(i_0)}$ and $X_1 \cap Z_{i_1}^{H(i_1)}$ are finite, therefore $(X_0 \cup X_1) \cap ( Z_{i_0}^{H(i_0)} \cap Z_{i_1}^{H(i_1)})$ is finite, contradicting our assumption that $X_0 \cup X_1, Z_{i_0}^{H(i_0)}, Z_{i_1}^{H(i_1)} \in \A$ and Claim 3.
% \smallskip

% \emph{Claim 5. $\A \setminus \P \neq \emptyset$.} Let $X = \{x_0, x_1, \dots \}$ defined as follows: let $x_n$ be the smallest element of $\bigcap_{k \leq n} Z_k^{H(k)}$ bigger than $2^{x_{n-1}}$ (or bigger than $0$ if $n = 0$). By construction $X$ cannot be an element of $\P$ as $|X \cap [x,2^x)| \leq 1$ for every $x$, and $X \in \A$ as $X \cap Z_k^{H(k)}$ is infinite for every $k \in \NN$.
% \smallskip


% \emph{Claim 6. $\P \cap \A$ is partition regular.} Clearly, $\NN \in \P \cap \A$. Fix $X_0,X_1$ such that $X_0 \cup X_1 \in \P \cap \A$. If $X_0, X_1 \in \P$, we are done as $\A$ is partition regular (Claim 4). Therefore, assume without loss of generality that $X_0 \notin \P$, and thus $X_1 \in \P$ by partition regularity of $\P$. Remains to show that $X_1 \in \A$. It is sufficient to show that $X_1 \cap Z$ is infinite for every $Z \in \A \cap \M$. Let $Z \in \A \cap \M$ (thus $Z \in \P$), $X_0 \cap Z$ is not in $\P$, therefore $X_1 \cap Z$ is in $\P$ by partition regularity of $\P$ (Claim 1), and $X_1 \cap Z$ is infinite.
% \smallskip

% By construction, $\A$ is $\M$-cohesive. Suppose for the contradiction that $\A$ is $\M$-minimal. Since $\P$ can be written as an intersection of $\Sigma^0_1$ classes, then by $\M$-minimality of~$\A$, $\A \subseteq \P$. This contradicts Claim 5.
\end{proof}

\end{document}
